\section{Introduction}
Let $k$ be a field. Recall that the \textit{Specht modules} for the symmetric group $S_n$ (over $k$) are a family of $kS_n$-modules
(we shall use right modules in this article) which are indexed by the partitions of $n$. We shall write the Specht module for $S_n$ which
is indexed by the partition $\lambda$ as $S^\lambda$. These Specht modules have a close relationship with the simple modules of
$kS_n$, and indeed if $kS_n$ is semisimple then the Specht modules are exactly the simple modules. Because of this and other properties,
the Specht modules for $kS_n$ have been the subject of intense study for decades, and a large and varied literature has built up around
them.

In this article we consider a class of modules for the \textit{wreath product} $S_m \wr S_n$ of two symmetric groups which are
analogous to the Specht modules for the symmetric group. These modules may be obtained from Specht modules for $S_m$ and $S_n$ via
a well-known method of constructing modules for wreath products, see for example \cite{CHTAN} or \cite[chapter 4]{JAMKER}.
We may alternatively obtain them as the cell modules of a certain cellular structure (in the sense of Graham and Lehrer) on the group
algebra $k(S_m \wr S_n)$. This cellularity was originally proved in \cite{GEGO}, while an alternative proof via the method of
\textit{iterated inflation} was given in \cite{GCSWPA}. The characterisation of these modules as cell modules
allows us to see at once that they bear exactly the same relation to the simple modules for the wreath product as the symmetric group
Specht modules bear to the simple modules for that group, and this justifies the name ``Specht modules''. Although the construction by
which these modules may be obtained is well-known, the author is not aware that these modules have previously been studied in the
literature as wreath product analogues of the symmetric group Specht modules.

A key fact in the theory of Specht modules is the result of James which we shall call the ``Specht branching rule'',
which gives a Specht filtration for the restriction of a Specht module from $kS_n$ to $kS_{n-1}$ with an
elegant combinatorial description of the set of Specht modules occurring in this filtration. Moreover, these multiplicities are
independent of the field $k$. The main results presented in this paper are two Specht branching rules for the wreath product of two
symmetric groups: the first describes the restriction of a wreath product Specht module from $k(S_m \wr S_n)$ to $k(S_{m-1} \wr S_n)$,
while the second describes the restriction to $k(S_m \wr S_{n-1})$. In both cases, we obtain a Specht module filtration with
multiplicities that do not depend on the field and which moreover have nice combinatorial descriptions.

Note that we are using the name ``branching rule'' in what is perhaps a slightly non-standard way. Indeed, in general group
representation theory, if we have some nested family of finite groups $G_1 \leq G_2 \leq \cdots \leq G_n \leq G_{n+1} \leq \cdots$, then a
\textit{branching rule} is a result describing the simple composition factors of the restriction of a simple module from $G_{n+1}$ to
$G_n$. Such branching rules are known in general for groups of the form $G \wr S_n$ for $G$ a finite group, see in particular
\cite{TSUWRBR} for a modular branching rule for $G \wr S_n$.

However, the results we present here are concerned specifically with \textit{Specht} modules, which do not always coincide with the simple
modules, and thus the results of this paper are not consequences of the known general results. However, our results are of a very similar
nature, and indeed our Specht modules are in fact the simple modules when the group algebra is semisimple (in both the symmetric group
and the wreath product case) and so in the semisimple case our Specht branching rules are in fact branching rules in the more usual
sense (and thus coincide with the known general results). The proofs of our results are however rather different to those of the known
general branching rules, since by specialising to the symmetric groups and focussing on the Specht modules, we can use combinatorial
methods rather than needing to appeal to, for example, Clifford theory as in the general case.


\section{Background}
We shall work over a field $k$ in this article. We shall often need to deal with tensor products of $k$-vector spaces, and we shall
abbreviate $\otimes_k$ to $\otimes$.

If $G$ is a group and $k$ is a field, then we shall write $kG$ for the group algebra of $G$ over $k$. By a \textit{$kG$-module}, we shall
mean a right $kG$-module of finite $k$-dimension.
If $G$ is a group with a subgroup $H$, then for a field $k$ we shall write the operations of induction and restriction of modules
between the group algebras $kG$ and $kH$ as $\uparrow^{G}_{H}$ and $\downarrow^{G}_{H}$, with the field being implicit.

Mackey's theorem is a fundamental result in finite group theory which describes the interaction of the operations of induction and
restriction. If $G$ is a finite group and $H$ is a subgroup of $G$, then for $g \in G$ we define $H^g$ to be the subgroup
$\{g^{-1}hg \mid h \in H\}$ of $G$, and we call this the \textit{conjugate subgroup of $H$ by $g$} (note that $H^g$ is isomorphic to $H$).
Further, if $X$ is a $kH$-module, then we define $X^g$ to be the $kH^g$-module with underlying vector space $X$ and action given by
$x(g^{-1}hg) = xh$ for $x \in X$ and $h \in H$. We call this the \textit{conjugate module of $X$ by $g$}.

\begin{theo}\label{mackey:thm}\textup{(}Mackey's Theorem \cite[Theorem 3.3.4]{BRCI}\textup{)}
Let $G$ be a finite group with subgroups $H$ and $K$, let $\mathcal{U}$ be a complete non-redundant system of $(H,K)$-double coset
representatives in $G$, and let $X$ be a right $kH$-module. Then we have a decomposition of right $kK$-modules
\[
X\!\uparrow^{G}_{H}\downarrow^{G}_{K} \;\;\cong\;
\bigoplus_{u \in \mathcal{U}} X^u\downarrow^{H^u}_{H^u\cap K}\uparrow^{K}_{H^u\cap K}.
\]
\end{theo}

\subsection{Filtrations and series}\label{filt_ser:subsec}
Let $kG$ be a group algebra over a field, let $M$ be a $kG$-module and $X_1,\ldots,X_t$ also be $kG$-modules. A \textit{series for $M$
over $X_1,\ldots,X_t$} is a series of submodules
\begin{equation}\label{M_series:eq}
  M = M_n \supseteq M_{n-1} \supseteq M_{n-2} \supseteq \cdots \supseteq M_1 \supseteq M_0 = 0
\end{equation}
such that each quotient $\frac{M_l}{M_{l-1}}$ is isomorphic to some $X_i$. If such a series exists, we say that $M$ has a
\textit{filtration by the modules $X_1,\ldots,X_t$}. If $\frac{M_1}{M_0} = M_1$ is isomorphic to $X_l$, then we say that $X_l$ occurs at
the \textit{bottom} of the filtration or series. Further, let $f:\{1,\ldots,n\} \longrightarrow \{1,\ldots,t\}$ be any function such
that for each $l$, $X_{f(l)}$ is isomorphic to the quotient of $M_l$ by $M_{l-1}$, and let $\alpha_i = |f^{-1}(i)|$. We then say that
\eqref{M_series:eq}
is a \textit{series for $M$ over $X_1,\ldots,X_t$ admitting the multiplicities $(\alpha_i)_i$}, and that $M$ has a
\textit{filtration by the modules $X_1,\ldots,X_t$ admitting the multiplicities $(\alpha_i)_i$}.
Note that we have not assumed that the modules $X_i$ are pairwise non-isomorphic. If there are isomorphisms between the modules
$X_i$, then the multiplicities in a filtration are \textit{not uniquely determined} by the series of submodules, and so the same series
of submodules will admit different multiplicities. Even if the modules $X_i$ are pairwise non-isomorphic, so that the multiplicities
\textit{are} uniquely determined by the series, the multiplicities are \textit{not} in general uniquely determined by the module,
as the same module can have two series of submodules where the $X_i$ occur with different multiplicities.

\subsection{Combinatorics}
We now review a few combinatorial concepts. We assume that the reader is already familiar with these notions, and so our treatment will
be brief.

Recall that a \textit{composition} of $n$ is a tuple of non-negative integers adding up to $n$. We call $n$ the \textit{size} of $\alpha$
and write $n = |\alpha|$. We call the elements of a composition its \textit{parts}, and the number of parts in a composition is its
\textit{length}. We shall adopt the common shorthand of using exponent
notation for repeated parts in a composition, so that for example we might write $(3,2^2,1^3)$ for $(3,2,2,1,1,1)$.
A \textit{partition} of $n$ is a composition of $n$ whose parts are all positive and appear in non-increasing order. We shall write
$\lambda \vdash n$ to mean that $\lambda$ is a partition of $n$. We note that the \textit{empty partition} $()$ is the unique partition
of size 0, and also the unique partition of length 0.

A simple total order on the partitions of an integer $n$ is the \textit{lexicographic order}, in which partitions are sorted by the
size of their first part, then by the size of their second part, and so on. Thus in this order, $(n)$ is the largest and $(1^n)$ the
least partition.
A more sophisticated (partial) order in the partitions of $n$ is the \textit{dominance order}, where $\lambda \trianglerighteq \mu$ if and
only if the sum of the first $t$ parts of $\lambda$ equals or exceeds the sum of the first $t$ parts of $\mu$ for all $t$.

The \textit{Young diagram} of a composition $\alpha$ is an arrangement of rows of boxes with a number of boxes on the $i$\tss{th} row
(counting downward) equal to the $i$\tss{th} part of $\alpha$.
If $\alpha,\gamma$ are compositions of $n$, then a \textit{tableau of shape $\alpha$ and type $\gamma$} is a Young diagram of shape
$\alpha$ where each box contains a positive integer $i$ such that for each $i \in \{1,\ldots,t\}$ where $t$ is the length of $\gamma$,
$i$ occurs exactly $\gamma_i$ times.
Note that since we allow zero parts in compositions, a Young diagram or tableau can have empty rows.

Now if $\lambda,\alpha$ are compositions such that $|\alpha|\leq|\lambda|$ and the Young diagram of $\alpha$ lies wholly inside the Young
diagram of $\lambda$ (\ie the length of $\alpha$ is at most the length of $\lambda$ and $\alpha_i \leq \lambda_i$ for all $i$ from 1 to
the length of $\alpha$), then for $\gamma$ a composition of $|\lambda|-|\alpha|$, we define a \textit{skew tableau of shape $\lambda
\setminus \alpha$ and type $\gamma$} to be a diagram obtained by removing the boxes of the Young diagram of $\alpha$ from $\lambda$ and
then filling the remaining boxes with positive integers such that each $i$ occurs $\gamma_i$ times, as for non-skew tableaux. Note that
the boxes of a skew tableau may be non-contiguous, as in the example below.

If a tableau (skew or non-skew) has its entries strictly increasing down each column and weakly increasing from left to right across each
row, we say that it is \textit{semistandard} (note that there may be gaps in the columns of a skew tableau). Thus for example
